% A mathematics paper in the AMS article class (amsart).
% The \address is printed at the end of the paper.
\documentclass[11pt]{amsart}
\usepackage{amssymb}

% Theorem-like environments, numbered within sections (1.1, 1.2, ...).
\newtheorem{theorem}{Theorem}[section]
\newtheorem{lemma}[theorem]{Lemma}
\newtheorem{corollary}[theorem]{Corollary}
\theoremstyle{definition}
\newtheorem{definition}[theorem]{Definition}
\theoremstyle{remark}
\newtheorem*{remark}{Remark}

\newcommand{\N}{\mathbb{N}}

\title{On the Infinitude of Primes}
\author{LeafTeX}
\address{Department of Mathematics, Example University}

\begin{document}

\maketitle

\begin{abstract}
We give a short proof that there are infinitely many primes and record
two simple consequences for the sequence of primes.
\end{abstract}

\section{Introduction}
For $x \ge 1$ let $\pi(x)$ be the number of primes $p \le x$. The classical
reference is Hardy and Wright~\cite{hw}.

\begin{definition}
A natural number $p > 1$ is \emph{prime} if its only positive divisors are
$1$ and $p$.
\end{definition}

\section{Main result}

\begin{theorem}[Euclid]
\label{thm:infinite}
There are infinitely many primes.
\end{theorem}

\begin{proof}
Suppose $p_1, \dots, p_n$ were all the primes, and let
\[
  N = p_1 p_2 \cdots p_n + 1.
\]
No $p_i$ divides $N$, yet $N > 1$ has a prime factor, a contradiction.
\end{proof}

\begin{corollary}
For every $n \in \N$ there is a prime $p > n$.
\end{corollary}

\begin{remark}
The argument of Theorem~\ref{thm:infinite} also gives $p_{n+1} \le p_1 \cdots p_n + 1$.
\end{remark}

\begin{thebibliography}{9}
\bibitem{hw}
G.~H. Hardy and E.~M. Wright, \emph{An Introduction to the Theory of
Numbers}, 6th ed., Oxford University Press, 2008.
\end{thebibliography}

\end{document}
